\section{Evidence-Centric Robustness: Noisy Retrieval and Distraction}

The first branch of robust RAG studies what enters the context. Early work on irrelevant context established that retrieval errors do not merely lower recall; they can actively redirect generation \citep{irrelevantContext2023}. Noise-aware training and robust RALM work then separated random negatives, hard negatives, counterfactual passages, and noisy retrieved evidence as different stressors rather than one generic perturbation \citep{raat2024,rrag037,rrag091}. Benchmarks such as RGB and URAG made this operational, testing negative rejection, integration, counterfactual evidence, and uncertainty under retrieval noise \citep{rgb2023,rrag001}. These works show that ``noise robustness'' is not a single axis but a family of distinct stressors.

Distraction work focuses on negative marginal utility: a passage is harmful not simply because it is irrelevant, but because it changes the answer, confidence, attribution, or reasoning path. Long-context and context-utility studies show that more context can increase the probability of including plausible but misleading passages \citep{rrag003,rrag039,rrag032}. Utility-oriented retrieval and downstream-aware metrics therefore ask whether a passage helps the generator rather than whether it only matches the query \citep{rrag058,rrag067,rrag073,rrag081,rrag088}. Context compression and filtering methods try to remove irrelevant spans or documents while preserving evidence needed for multi-hop reasoning \citep{rrag044,chainOfNote2023,sufficientContext2024,rrag043,rrag048,rrag055}.

The later evidence-centric literature broadens noise beyond the top-$k$ ranking to query phrasing \citep{queryRobustness2025}, context ordering and length \citep{rrag140,rrag143,rrag154,rrag166,rrag168,rrag170}, and pipeline-level construction, compression, and retrieval defects \citep{rrag174,rrag179,rrag180,rrag181,rrag182,rrag183,rrag184}, with recent systems and benchmarks adding stress tests for retrieval defects, context utility, and robust fine-tuning \citep{rrag200,rrag201,rrag206,rrag210,rrag212,rrag215,rrag218,rrag219,rrag224,rrag225,rrag226}. The common lesson is that retrieval should be evaluated as context construction, not recall.

Venue-accepted work reinforces the trend---efficient RAG and reranking \citep{venueAcc013,venueAcc017,venueAcc018}, graph- and knowledge-guided retrieval \citep{venueAcc004,venueAcc006,venueAcc008,venueAcc009,venueAcc012,venueAcc023}, long-context benchmarks \citep{venueAcc028,venueAcc030}, and compression and retrieval-utility methods \citep{venueExtra003,venueExtra015,venueExtra007,venueExtra014,venueExtra002,venueExtra009,venueExtra013,venueExtra012,venueExtra008}---all treating retrieval as evidence construction rather than recall.

\section{Model-Centric Robustness: Conflict, Faithfulness, and Attribution}

Once evidence enters the prompt, robustness depends on how the generator uses it. Knowledge-conflict papers study contradictions between parametric memory and retrieved evidence, or among retrieved sources \citep{tugOfWar2024,faithfulRAG2025,astuteRAG2025,rrag046}. These papers show that a model can over-trust retrieved evidence, resist correct retrieval because of parametric priors, or silently choose one source without exposing disagreement. Conflict-aware methods therefore localize discrepancies, compare internal and external knowledge, or introduce source-aware reasoning before generation \citep{rrag008,rrag013,rrag017,rrag024,rrag029,rrag031,rrag036,rrag038}.

Faithfulness and attribution papers study whether generated claims are supported by evidence. RAGTruth and related grounding work provide fine-grained hallucination or support labels, while source-preference and entity-conflict studies examine whether models cite and trust the right evidence \citep{ragTruth2024,rrag110,rrag112,rrag115,rrag132}. Document-level citation is often too coarse, so claim- or sentence-level attribution is more informative but costlier to annotate \citep{rrag076,rrag077,rrag082,rrag086,rrag089,rrag092,rrag100}.

The broader map extends this to source reliability, contextual faithfulness, real-document conflicts, and hidden-state or logit interventions \citep{rrag118,rrag130,rrag147,rrag150,rrag152,rrag163,rrag176,rrag187}, and to richer benchmarks, source-aware generation, and practical pipelines \citep{rrag202,rrag204,rrag207,rrag209,rrag214,rrag216,rrag220,rrag222,rrag223,rrag227}. One distinction recurs: evidence can be present but ignored, cited but unsupported, or correct but less trusted than a stale parametric belief.

Venue-accepted faithfulness work adds the same distinction: grounded attribution and refusal \citep{venueAcc001}, context-faithfulness benchmarks under misleading evidence \citep{rrag040}, hallucination and mechanistic detectors \citep{venueAcc020,venueAcc026}, and parametric--contextual knowledge reconciliation \citep{venueAcc027} together show that robust RAG needs both evidence-use diagnostics and source-reliability judgments.

\section{Adversarial Evidence Corruption and Security}

Adversarial RAG work asks what happens when evidence unreliability is created by an adaptive attacker. Corpus and retrieval poisoning papers show that small numbers of malicious passages can manipulate dense retrieval or steer final answers \citep{poisoningRetrievalCorpora2023,poisonedRAG2024,badRAG2024,rrag049,rrag057}. Certified and aggregation-based defenses study robustness under bounded corruption assumptions \citep{certifiedRAG2024,rrag016,rrag020}, while reranking and detection defenses use gradient, attention, clustering, or anomaly signals to demote suspicious evidence \citep{rrag004,rrag009,rrag011,rrag015,rrag017}.

Prompt-injection and backdoor work targets a different interface. Instead of inserting false facts, the attacker may insert instructions that the model treats as commands, or triggers that activate malicious behavior through retrieved contexts \citep{aip2025,rrag005,trojanRAG2024,phantom2024,rrag010,rrag018}. These attacks require retrieval realism: a malicious document must be retrieved under plausible queries, not only under attacker-crafted prompts. This is why we separate factual poisoning, indirect prompt injection, backdoor triggers, and prompt-template attacks as different threat games even when they all appear as retrieved text.

The security map is broad and fast-moving, spanning poisoning benchmarks and arenas \citep{rrag019,rrag023,rrag033,rrag035}, traceback and forensics \citep{rrag014,rrag128}, privacy and membership inference \citep{rrag060,rrag063,rrag096,rrag103}, opinion or fairness manipulation \citep{rrag104,rrag162}, graph, multimodal, or reasoning-style poisoning surfaces \citep{rrag051,rrag127,rrag175}, and many preprint defenses and attacks \citep{rrag022,rrag025,rrag047,rrag052,rrag068,rrag097,rrag114,rrag126,rrag149,rrag177,rrag197,rrag221}. Its use is not to rank attacks but to motivate capability-normalized reporting.

Venue-accepted security work extends the surface to agent and memory poisoning, data extraction, medical-domain risk, robust web search, and semantic conflict-free filtering \citep{rrag053,rrag090,rrag021,rrag017,venueAcc014}, comparable only once attacker control over corpus, prompt, query, memory, or retrieved evidence is specified.

\section{System Control: Calibration, Abstention, and Adaptive Retrieval}

System-control work converts evidence judgments into actions. Answerability and sufficient-context papers ask whether the retrieved context contains enough evidence to answer \citep{sufficientContext2024,rrag111,rrag141,rrag193}. Uncertainty and calibration work asks whether confidence reflects answer correctness, evidence sufficiency, retrieval usefulness, or selective risk \citep{calibRAG2025,uncertaintyRAG2024,rrag001,rrag007,rrag012,rrag093}. Adaptive retrieval work asks when the system should retrieve, retrieve again, or answer from internal knowledge \citep{selfRAG2023,whenToRetrieve2024,crag2024,ctrla2025,rrag034}.

The control literature highlights a recurring failure: confidence is not the same as evidence sufficiency. A model can be confident in an unsupported answer, or uncertain even when evidence is present but hard to integrate. Trust-or-abstain and self-aware RAG methods therefore model whether to answer, abstain, retrieve more, or ask for clarification \citep{rrag002,rrag006,rrag041,rrag064,rrag066,rrag074,rrag075,rrag087,rrag095,rrag105,rrag108}. Multi-turn and underspecified-query benchmarks add another layer: the correct action may be clarification before retrieval, not direct answering \citep{rrag141,rrag146,rrag153}.

Recent work extends control to cost-aware retrieval, conformal or set-valued prediction, risk-sensitive abstention, and adaptive policies over multiple retrieval steps \citep{rrag159,rrag164,rrag167,rrag169,rrag173,rrag178,rrag186,rrag191,rrag192,rrag198,rrag205,rrag228}. These papers motivate an action-control view: robust RAG is not only about reducing errors, but about choosing the right response under uncertainty and cost.

Venue-accepted control work makes the action view explicit: calibration framed as retrieval-and-answer control \citep{venueAcc003}, learned adaptive retrieval and critique \citep{venueAcc016}, query routing across systems or long-context processing \citep{rrag039,venueAcc028,venueAcc030}, multi-step chain-of-retrieval \citep{venueAcc010,venueExtra001}, and feedback- or RL-driven policies for when to retrieve, rewrite, and answer \citep{venueExtra010,venueExtra011,venueExtra016,venueExtra006}.

\section{Evaluation and Benchmarks}

Evaluation papers provide the diagnostic infrastructure for robust RAG. RGB decomposes RAG reliability into multiple abilities \citep{rgb2023}; RAGTruth provides fine-grained hallucination annotations \citep{ragTruth2024}; query-robustness and context-utility work stress retrieval interfaces \citep{queryRobustness2025,rrag058}; security benchmarks and arenas evaluate poisoning and defense assumptions \citep{rrag019,rrag023,rrag035}; and uncertainty benchmarks evaluate whether systems know when their answers are reliable \citep{rrag001,rrag093}. Survey papers on RAG, trustworthy RAG, secure RAG, and evaluation provide the broader landscape \citep{ragSurvey2024,ragComprehensive2025,ragEvaluationSurvey2024,trustRAGCompass2024,trustworthyRAGSurvey2025,securityRAGSurvey2026,rrag134,rrag151}.

Accepted benchmark work adds grounding and hallucination verification \citep{venueAcc015,venueAcc020,venueAcc026} and multimodal or graph RAG benchmarks \citep{venueAcc002,venueAcc004,venueAcc006,venueAcc019}, treated as boundary evidence rather than the empirical backbone for text RAG.

What is still missing is joint evaluation. Existing benchmarks often isolate one failure mode: noise, conflict, poisoning, or unanswerability. Real systems face combinations. A benchmark that includes a stale source, a hard negative, insufficient evidence, and a malicious document will test a different capability from any single-axis benchmark. This motivates the action-trace blueprint: robust-RAG evaluation should score evidence state, answer correctness, action choice, cost, and threat assumptions together.
